\chapter{Vector Analysis}
\renewcommand{\currentchapterpath}{chapters/ch02/}

\begin{flushright}
\textit{“引其纲，万目皆张。” ——《吕氏春秋·用民》}
\footnote{From the \textit{Lüshi Chunqiu}: 
``Pull the main cord of a fishing net, and every mesh opens.'' 
The net's main cord represents the key, fundamental, 
or governing element of a matter: once it is understood, 
its many connected parts become clear and can be put in order.}
\end{flushright}


A \textit{vector} is a mathematical object that has both magnitude and direction,
which distinguishes it from a \textit{scalar}, which has only magnitude.
It is widely used in physics and engineering to describe quantities such as 
force, velocity, and acceleration in mechanics, and 
electromagnetic fields in electromagnetism, and fluid flow in fluid dynamics.
Its dimension can be more than three, 
extending to a more general higher-dimensional spaces. The study of 
vector and vector spaces helps in understanding 
more abstract concepts in physics which shall be encountered in later studies.


\chapterinput{basics}
\chapterinput{differentials}
\chapterinput{index_notation}
\chapterinput{tensors}
\chapterinput{exercises}